%!TEX root = ../template.tex %

% !BIB references = ../Bibliography/bibliography.bib

% !TeX spellcheck = en_US
%%%%%%%%%%%%%%%%%%%%%%%%%%%%%%%%%%%%%%%%%%%%%%%%%%%%%%%%%%%%%%%%%%%%
%% chapter3.tex
%% NOVA thesis document file
%%
%% Chapter with a short latex tutorial and examples
%%%%%%%%%%%%%%%%%%%%%%%%%%%%%%%%%%%%%%%%%%%%%%%%%%%%%%%%%%%%%%%%%%%%

\typeout{NT FILE chapter3.tex}%

\chapter{Quantum Mechanics of Electrons in Crystalline Solids}
\label{ch:quantum_mechanics_solids}

We have seen in the previous chapter, that even though classical models explain fairly well many phenomena in electrodynamics, both qualitatively and quantitatively, classical physics cannot predict the precise values of their respective parameters, nor can it explain why they work so well for certain materials, while failing to describe others.
To have a deeper understanding of these models and their regimes of validity, or to faithfully predict material properties in general, one has to adopt a microscopic description of the systems we want to study.
This means that we have to look at the chemical composition of the material and at the spatial arrangement of the atoms making up the material.
One has to look at the finer physics at length scales of the order of the angstrom, where classical physics is insufficient.

The tools that allow us to study systems at this level are provided by quantum mechanics, whose father is Max Planck, the german physicist who solved the problem of the ultraviolet catastrophe---an ideal black body radiates an infinite amount of energy at higher frequencies deep into the ultraviolet range of the electromagnetic spectrum.
This was the classical prediction by the Rayleigh--Jeans law, which obviously makes no sense.
It was the hypothesis by Planck that radiation is emitted in discrete amounts of energy, called quanta, that solved the mystery of the black body, giving birth to a new field in physics, which catalyzed a scientific revolution in the previous century.
Quantum mechanics allowed to solve other problems in physics until then considered open, such as why does not the electron approach the nucleus of the atom as it orbits around it, since it should lose energy due to its curvilinear motion as predicted by the laws of classical electrodynamics.

With tools provided by the underlying principles and laws of quantum mechanics, along the years many techniques have been developed since its birth to attack many of varying degree of complexity.
We keep our discussion narrowed to solids with a periodic atomic structure, called crystals, until the end of this thesis, and leave all other types of systems (amorphous solids, such as glasses, crystalline solids with defects, few-particle systems, etc.) out of the discussion.

In the field of solid state physics, it is customary to separate the techniques in first-principles based, called \abinit~techniques and are usually computationally intensive, and second-principles based techniques, which are usually much less computationally demanding.
In this thesis \abinit~approaches will be discussed on a passing note, as those can easily materialize into a thesis in itself, but are nonetheless very relevant in this work.

We start the chapter by...

\section{The Hamiltonian of a Crystal}
\label{sec:introductory_remarks}

\subsection{Born-Oppenheimer Approximation}
\label{subsec:born_oppenheimer_approximation}

\subsection{Non-interacting limit}
\label{subsec:non_interacting_limit}

\section{Single-Particle Properties of Electrons in Crystals}
\label{sec:single_particle_properties}

\subsection{Bloch's Theorem}
\label{subsec:bloch_theorem}


\subsection{Electronic band structure (methods)}
\label{subsec:electronic_band_structure}

\todo{Insert here the text from the notes}

\subsubsection{Tight-binding method}
\label{subsubsec:tight_binding_method}

Mention briefly free-electron model for metals, and introduce \tb models as the opposite limit for semiconductors and insulators.

Add hBN minimal \tb model as a prime example relevant for later.

\subsubsection{\textit{Ab initio} methods}
\label{subsubsec:abinit_methods}

Briefly discuss \dft and Wannierization, as we need it for the framework of the screening and excitons with the code.

\section{Many-Particle Description of Solids}
\label{sec:many_particle_description}

\subsection{Many-body perturbation theory and $GW$ approximation}
\label{subsec:many_body_perturbation_theory}

We need this to introduce the \rpa and Dyson's equation for the screened potential and introduce the polarisation and its non-interacting version.

\section{Dielectric Function and Screening: Quantum Picture}
\label{sec:dielectric_function_quantum}

Brief introduction stressing the difference between full, \rpa and non-interacting polarizability and dielectric function within self-consistent framework~\cite{Cohen_Louie_2016} and ~\cite{Bruus_Flensberg_2004}.
Then

\subsection{RPA dielectric function}
\label{subsec:rpa_dielectric_function}

\subsection{RPA screened Coulomb potential}
\label{subsec:rpa_screened_coulomb_potential}



